% The ulthesis class and most of the following came from
% https://prancer.physics.louisville.edu/ulthesis/

\documentclass{ulthesis}
\author{Joseph Downs}
% Your dissertation title should match pdf title above
\title{OpenZL's Effectiveness in Compression of Scientific Data}

../tex-cfg/base-cfg.tex
../tex-cfg/programming-cfg.tex

\bibliography{refs.bib}

% Change figure and table caption label styles

\usepackage[labelsep=period, labelfont=bf]{caption} % boldface with period

\usepackage{chngcntr}   % use to serially number figures and tables 
\usepackage{relsize}    % for BABAR text logo
\usepackage{subcaption} % for side by side figures

%=============================== Custom Commands ===============================
\newcommand{\openzl}{OpenZL}
\newcommand{\metainc}{Meta Platforms,~Inc.}


\makeatletter
\renewcommand\@memmain@floats{%
% Enable serially numbered figures, tables and equations

\counterwithout{figure}{chapter}
\counterwithout{table}{chapter}
\counterwithout{equation}{chapter}}

% Suppress section numbering in TOC

% \renewcommand{\cftsectionpresnum}{\begin{lrbox}{\@tempboxa}}
% \renewcommand{\cftsectionaftersnum}{\end{lrbox}}
% \makeatletter
% \setlength{\cftsectionnumwidth}{0pt}

% Restarts section numbered from 1 on each chapter

% \renewcommand{\thesection}{\arabic{section}}

%%%%%%%%%%%%%%%%%%%%%%%%%%%%%%%%%%%%%%%%%%%%%%%%
\def\klong{\ensuremath{K_{\textrm L}^{0}} }

\def\babar{\mbox{\slshape B\kern-0.3em{\smaller[2] A}\kern-0.1em B\kern-0.3em{\smaller[2] A\kern-0.2em R}} }
\def\pt{\ensuremath{p_{\textrm T}} }
%%%%%%%%%%%%%%%%%%%%%%%%%%%%%%%%%%%%%%%%%%%%%%%%
\begin{document}

% Aliases used in bib file for journal titles
%\include{journals}

%================================= Front matter ================================

% Your current highest degree
\authordegree{B.S. in Computer Science and Engineering, 2025}

% This dissertation degree
\dissertationdegree{Master of Engineering}

% The dissertation discipline (may differ from the department name) for example "Physics"
\dissertationdiscipline{Computer Science and Engineering}

% Month and year of degree award
\dissertationdate{December 2026}

% Month day and year of dissertation approval
\approvaldate{}

% College or school
\school{Speed School of Engineering}

% Department in which the candidate is enrolled
\department{Department of Computer Science and Engineering}


% Advisors and committee members with appropriate honorifics
\advisor{Dr. Adrian Lauf}
\firstmember{Co-advisor or  first member}
\secondmember{Second member}
\thirdmember{Third member}
% \fourthmember{Fourth member}


% Use the same keywords for the pdf headers above
\keywords{Keywords or phrases separated by commas}


%This builds the title page and other front content
\frontmatter

\maketitle
% Now we add edited content page by page.
% Retain each section, edit as needed.
% Keep this order of sections and end with the tables of contents.

% \begin{dedication}

% \end{dedication}
% Insert external PDF page as needed by uncommenting the following
% \includepdf[pages=-]{Biswas_Electronic_Dissertation_Signature_Page_signed.pdf}

% An section with acknowledgments is conventional

\begin{acknowledgments}

\end{acknowledgments}


% Dissertation abstract should summarize the content

% The best practice is to not have citations in the abstract
% Avoid language which is specialized and acronyms that are not defined

\begin{dissertationabstract}
  As dataset sizes in AI/HPC workloads increase, data compression could become
useful to speed up overall execution time by reducing the time taken to transfer
between nodes on the network. However, this is not useful if the time
\emph{cost} of compression/decompression outweighs the time \emph{savings} in
data transmission over the network. This paper explores methods of compressing
data --- using OpenZL, a recent lossless compression framework --- in various
scenarios, i.e., compression before execution, compression during execution, and
no compression. I compare and analyze the time taken with various compression
algorithms and datasets. With the scenarios and datasets tested, \{all
scenarios; some scenarios; no scenarios\} exhibited speedup as a result of
compression with OpenZL.

\end{dissertationabstract}



% These commands will generate the required tables of contents. 
% The \clearpage commands insure that the next one begins on a new page.

\tableofcontents* \clearpage
\listoftables \clearpage
\listoffigures \clearpage


%================================= Main matter =================================

% Note we need this to insure proper formatting for the body of the dissertation.
\mainmatter

% Use separate files for chapters and include the chapter titles here. Having
% the titles here is necessary to get them properly indexed and formatted. The
% \input{file} will pull in the content of that file. You do not have to give an
% extension ".tex" but it anticipates the file has that.

\chapter{Introduction}
\label{sec:intro}
\section{Purpose}
High-performance computing dataset sizes are
increasing.\cite{reed_exascale_2015}. As \citeauthor{underwood_fraz_2020} state
in \cite{underwood_fraz_2020}, ``Today’s scientific research applications
produce volumes of data too large to be stored, transferred, and analyzed
efficiently because of limited storage space and potential bottlenecks in I/O
systems.'' \cite{foster_computing_2017} provides three examples of applications
generating large volumes of data:
\begin{itemize}
\item Climate science simulations checkpointing generating \qty{260}{\tera\byte}
  every 16~seconds.
\item Singe XGC (``X-point included Gyrokinetics Code'') ITER fusion models
  producing ``hundreds of petabytes''\cite{foster_computing_2017}.
\item Femtosecond-level material science simulations producing ``exabytes on
  exascale computers''\cite{foster_computing_2017}.
\end{itemize}

\cite{foster_computing_2017} notes this as a problem due to 2024-projected (at
the time) compute rate outperforming disk write rate by 50 times. Because I/O
throughput scales more slowly than computational speeds, ``We can no longer
output every piece of information that we might ever possibly
want''\cite{foster_computing_2017}.

\section{OpenZL}
The meat and potatoes of this thesis rests on OpenZL and its innovations. What
is OpenZL? OpenZL is a graph-based compressor released by \metainc{} in a
whitepaper, ``OpenZL: A Graph-Based Model for
Compression''\cite{collet_openzl_2025}, in 2025. It allows for greater
performance over generic compressors without the complexity and challenges of
application-specific compressors. The compressed data format is ``a
self-describing wire format, \ldots which can be decompressed by a universal
decoder''\cite{collet_openzl_2025}. Because of this, it allows compressors
tailored to datasets with universal decoding. The whitepaper lists the following
benefits:
\begin{itemize}
\item Minimal upfront investment
\item Flexibility of input data formats
\item Minimized security surface
\item Easy iteration
\item Fast deployment
\item Support for high-cardinality data
\end{itemize}


\chapter{Observations}
%\input{observations.tex}

\chapter{Analysis}
%\input{analysis.tex}

\chapter{Conclusions}
%\chapter{Conclusions}
\label{sec:label}



%================================= Back matter =================================

% The layout changes for the last part and you need this command
\backmatter

% You can change the style of the bibliography for whatever would be accepted by
% your discipline and the University format checkers.

% There are several bibliography managers that may be used with LaTeX. Here we
% have defaulted to bibtex, but if this does not work for you try others.

% Style "apalike" is similar to the Harvard style. If you prefer numbered
% references, try "plain".

% The ulthesis package includes a copy of the Optical Society of America's
% bibliography style which retains titles and provides a clean looking
% bibliography in order of citation in the body. It is used here by default.

% If you want a true Harvard style you can use "biblatex" or "natbib" insdotead of
% bibtex, which requires some other modifications. There is help for this on
% Sharelatex and on Overleaf, as well as through Google searches.

%\bibliographystyle{apalike}
%\bibliographystyle{plain}
%\bibliographystyle{osajnl}
%\bibliographystyle{ieee}
%\bibliography{dissertation}

\printbibliography

% Appendices and the Vita are treated as chapters in the back matter

\chapter{Appendix A: }

\chapter{Appendix B: }

\chapter{Curriculum Vitae}
%\input{vita.tex}


\end{document}
